%%%%%%%%%%%%%%%%%%%%%%%%%%%%%%%%%%%%%%%%%%%%%%%%%%%%%%%%%%%%%%%%%%%%%%%%%%%
%%  conclusion.tex  –  Conclusion (not a numbered chapter)
%%  v9 – Progressive summary of all defended results following the
%%       sequential chain: CT dynamics → distributed → efficiency →
%%       adaptation → certification → domain → platform.
%%       Matches the 5-chapter structure.
%%%%%%%%%%%%%%%%%%%%%%%%%%%%%%%%%%%%%%%%%%%%%%%%%%%%%%%%%%%%%%%%%%%%%%%%%%%

\chapter*{CONCLUSION}
\addcontentsline{toc}{chapter}{Conclusion}
\markboth{CONCLUSION}{CONCLUSION}

\section*{Brief Characterisation of the Main Results
(Eight Statements Aligned with the Eight Research Tasks)}

The main scientific results of the dissertation are systematised
below as eight statements in strict correspondence with the eight
research tasks set in the Introduction. Each statement records
how the corresponding task was solved and where the solution is
presented.

\begin{enumerate}[leftmargin=2em, itemsep=0.3em]

\item \textit{On the task ``Systematic review of the approaches
used in modern hybrid intrusion-detection and -prevention systems
and identification of their disadvantages and vulnerabilities'':}
Chapter~\ref{ch:literature} provides a systematic review
of approaches used in modern Hybrid~IDPS; documented critical
model failures in healthcare (AUC drop to 25\,\%), finance
($25.7\,\%\to 68.6\,\%$ errors), industry ($94.2\,\%\to 30\,\%$
accuracy) and cybersecurity ($-59\,\%$ detection), as well as a
robustness gap of 25--65 p.p.\ under adversarial influence.

\item \textit{On the task ``Review of methods to improve
robustness against adversarial attacks and to ensure computational
efficiency'':} five families of models and algorithms for
Hybrid~IDPS (statistical, classical ML, deep learning, hybrid,
game-theoretic; Fig.~\ref{fig:orga_models_idps}) are
systematised, and a comparative analysis of AI-augmented
commercial products (CrowdStrike Charlotte~AI, SentinelOne
Purple~AI, Microsoft Security Copilot, Google Mandiant+Gemini,
Cisco Hypershield, IBM QRadar AI, Palo~Alto Cortex~XSIAM) is
carried out, which records the absence of a unified architecture
simultaneously providing certified robustness, distributed
privacy, computational efficiency and calibrated accuracy on
dynamic graph structures.

\item \textit{On the task ``Formalise the mathematical model of
the hybrid intrusion-detection-and-prevention system and state the
problem of improving its robustness and efficiency'':}
Chapter~\ref{ch:methods} establishes a unified tuple-based
mathematical model of Hybrid~IDPS, gives a formal definition of
the six properties P1--P6 and of the twelve-cell
``condition\,---\,requirement'' matrix, and explicitly poses the
multi-criterion problem of improving robustness, accuracy,
efficiency and privacy.

\item \textit{On the task ``Develop algorithms and models for
improving the robustness and efficiency of the hybrid
intrusion-detection system under distributed and non-stationary
conditions'':} seven models M1--M7 are developed, providing
continuous-time graph dynamics with a Lipschitz/Gronwall
certificate (\textbf{CT-TGNN/SDE-TGNN}; 99.0\,\% F1 on
CIC-IoT-2023, certified radius $\rho=0.25$), Byzantine-resilient
federated learning with differential privacy (\textbf{FedLLM-API};
87.1\,\% F1 under 40\,\% corrupted clients, $\varepsilon=0.85$),
multi-level temporal representations (\textbf{TripleE-TGNN};
96.8\,\% F1 on multi-stage attacks), a state-space model with
poisoning resilience (\textbf{MambaShield}; $\mathcal{O}(L\log L)$,
91.4\,\% F1 under 25\,\% poisoning, 0.5~ms latency), the
stochastic transformer with UC-HGP (ECE\,=\,0.078;
$-43\,\%$ SOC workload), and the unified game-theoretic
certification framework (\textbf{M7-Certification}; 94.7\,\% F1
against adaptive adversaries).

\item \textit{On the task ``Develop software implementing the
proposed models and algorithms (formulating requirements, use
cases, architecture and program realisation)'':}
Chapter~\ref{ch:platform} introduces the software platform
\textbf{RobustIDPS.ai} (DOI~10.5281/zenodo.19129512) realising
Models M1--M7 within a single system: 17 neural-network models,
10 functional groups, 51 pages of a Progressive Web Application
(PWA), a four-level LLM defence (CyberSecLLM, 94.5\,\% detection
across 517 LLM attack scenarios, 89.1\,\% zero-shot F1 on the
cybersecurity benchmark, $+20.5$ p.p.\ over GPT-4); functional
requirements, use cases and system architecture are formulated
(FastAPI/PyTorch~2.x/PyTorch~Geometric, React~18/PWA,
PostgreSQL~16/Redis~7, Docker~Compose; four LLM providers and a
Suricata EVE-JSON plug-in; see Section~\ref{sec:tech_stack_ui}).

\item \textit{On the task ``Conduct experimental research of the
developed algorithms and models on benchmark datasets'':}
Chapter~\ref{ch:experiments} carries out a comprehensive
experimental programme on six network-traffic benchmark datasets
(CIC-IoT-2023, CSE-CIC-IDS2018, UNSW-NB15, MS-OUEDA, Container
Security, Edge-IoT), covering 84.2~million labelled samples,
84 attack types, and 23 metrics of quality, robustness,
efficiency, privacy and calibration; all 12 cells of the
``condition\,---\,requirement'' matrix are confirmed with a
superiority over baseline methods of 2.9--5.6 p.p.

\item \textit{On the task ``Integrate the developed software into
a real network-traffic-analysis pipeline'':} \textbf{two
implementation acts} have been received that confirm integration
of the developed RobustIDPS.ai software into real production
pipelines for network-traffic analysis and infrastructure
protection: from NII~IKT~LLC (Novosibirsk, 03~June 2026;
integration of poisoning-resilient models into the operating
intrusion-detection system and active use of the RobustIDPS.ai
platform by in-house cybersecurity specialists) and from
Ksyrix~LLC (Moscow, 2026; integration into the existing
network-traffic-analysis and infrastructure-protection software~---
multi-stage attack detection, identification of previously unknown
patterns, and adaptation to drift).

\item \textit{On the task ``Demonstrate that the developed
software admits reproduction in another practical area (on the
example of UAV attacks)'':} Chapter~\ref{ch:uav} implements the
three-tier ``onboard\,---\,droneport\,---\,cloud'' architecture
applying M1--M7 to the protection of computer-vision onboard
systems of unmanned aerial vehicles: the onboard tier runs the
MambaShield inference with per-flow latency 0.5~ms, the droneport
aggregates updates via FedLLM-API, and the cloud tier realises
the game-theoretic M7 certification. Domain independence of the
developed mathematical and algorithmic framework is confirmed;
the third implementation act (AI Center of NRNU MEPhI, Moscow) is
in the course of formalisation and will document this
reproducibility.

\end{enumerate}

\noindent An expanded discussion of the results listed above is
given in the following sections of this Conclusion.

\bigskip

This dissertation addressed the fundamental challenge of ensuring
the robustness, efficiency, accuracy, and privacy of \emph{Hybrid
Intrusion Detection and Prevention Systems} (Hybrid~IDPS)
operating under adversarial, distributed, and non-stationary
conditions. The research was motivated by documented failures of
existing IDPS solutions in healthcare, finance, industrial IoT,
autonomous systems, and cybersecurity, where system robustness was
shown to lag detection accuracy by 25--65~percentage points. A
unified mathematical model of Hybrid~IDPS was formalised as the
tuple $\mathcal{H} = (\mathcal{V}_t, \mathcal{E}_t, X_t, A_t,
f_\theta, g_\psi, \pi)$, and six cross-cutting properties of the
system (adversarial resilience, robustness, efficiency, distributed
environment, the human factor, and non-stationarity) were
systematically decomposed into seven sub-problems addressed by seven
novel Models and algorithms~M1--M7. Graph neural networks are used
in this work as one branch of the family of models applicable to
Hybrid~IDPS, alongside statistical, classical~ML, hybrid, and
game-theoretic approaches. Consistent with the title of the
dissertation, the contribution is positioned as the
\emph{development of adversarial resilient models based on
artificial intelligence in hybrid intrusion detection and
prevention systems for network security}, with the
network-traffic protection sector (NIDPS) serving as the primary
deployment instance and the UAV onboard sector as the
domain-transfer confirmation. The following results were obtained.


%%---------------------------------------------------------------------
\section*{Principal Results}
%%---------------------------------------------------------------------

\noindent\textbf{1. CT-TGNN and SDE-TGNN: Certified Continuous-Time Graph Dynamics (Gap~1).}
The first continuous-time temporal graph neural network was developed, unifying graph message passing with neural ODE dynamics on evolving topologies. Adjoint sensitivity training achieves $O(1)$ memory; Lipschitz-bounded dynamics yield certified robustness through the Gr\"{o}nwall inequality. The stochastic extension (SDE-TGNN) replaces deterministic dynamics with It\^{o} SDEs and derives analytical moment propagation from the Fokker-Planck equation. Experimental validation on three benchmarks totalling 52 million samples demonstrates 99.0\% average weighted F1, expected calibration error of 0.042, certified adversarial robustness of 76.2\% at perturbation radius $R=0.25$, and cross-domain transfer of 90.5\% F1 without fine-tuning.

\noindent\textbf{2. FedLLM-API: Privacy-Preserving Distributed Graph Learning (Gap~2).}
An original Byzantine-resilient federated graph optimiser was developed, combining coordinate-wise trimmed-mean aggregation with $(\varepsilon{=}0.85,\,\delta{=}10^{-5})$ differential privacy and Wasserstein optimal-transport alignment. Convergence at $\mathcal{O}(1/\sqrt{T})$ is proven under 30\% corruption. With 40\% Byzantine participants, FedLLM-API achieves 87.1\% accuracy versus 61.4\% for FedAvg. LLM integration enables 87.6\% F1 on zero-shot detection of novel attack categories.

\noindent\textbf{3. TripleE-TGNN: Multi-Level Temporal Graph Representation (Gap~3).}
A novel triple-embedding temporal graph architecture was developed, learning service-, trace-, and node-level representations via dual attention with elastic weight consolidation against catastrophic forgetting. Overall accuracy of 96.8\% with 40\% computational cost reduction, maintaining 94.7\% accuracy without retraining under continuous topology evolution.

\noindent\textbf{4. MambaShield: Efficient State-Space Inference (Gap~4).}
A new selective state-space architecture was developed, reducing inference complexity from $\mathcal{O}(L^{2})$ to $\mathcal{O}(L\log L)$ via FFT convolutions with progressive adversarial distillation. Under 25\% training-time poisoning, accuracy is 91.4\% versus 73.8\% for undefended training. PAC-Bayesian generalisation certificates confirm the bound with 2.7\% margin. Inference latency of 0.5\,ms per flow is the lowest of all models.

\noindent\textbf{5. Stochastic Transformer and UC-HGP: Uncertainty-Calibrated Adaptation (Gap~6).}
The first uncertainty-calibrated GNN inference combining variational Bayesian attention with hierarchical Gaussian processes was developed, achieving epistemic--aleatoric decomposition with ECE of 0.078 (59\% reduction over deterministic attention) and 43\% analyst workload reduction through uncertainty-driven triage.

\noindent\textbf{6. Game-Theoretic Certification: Unified Robustness Framework (Gap~7).}
An original robustness certification framework was constructed via Stackelberg equilibria with no-regret multiplicative weights achieving $R(T)\leq 2\sqrt{T\log|\mathcal{C}|}$ regret. Consistency regularisation couples all seven methods with proven convergence. Accuracy of 94.7\% against adaptive adversaries versus 87.6\% for non-strategic approaches.

\noindent\textbf{7. CyberSecLLM: Domain-Specific Foundation Model (Gap~5).}
A novel domain-specific foundation model with ODE-augmented selective state-space blocks was developed, achieving 89.1\% average zero-shot accuracy on the 11-task CyberSecBench evaluation --- a 20.5-point improvement over GPT-4 --- while processing million-token graph event sequences in 2.3 seconds.

\noindent\textbf{8. Domain Adaptation.}
All methods were instantiated in the cybersecurity domain through a formal data-to-graph transformation pipeline, a unified 34-class attack taxonomy, an 83-feature engineering pipeline, and domain-specific hyperparameter configurations. The adaptation demonstrates the portability of the mathematical framework to a concrete deployment environment.

\noindent\textbf{9. Experimental Validation.}
Comprehensive experiments across six benchmark datasets totalling 84.2 million labelled samples, 84 attack categories, and 23 evaluation metrics confirm that the unified framework outperforms the best baselines by 2.9--5.6 percentage points. Ablation studies verify that every component contributes measurably. Adversarial robustness is demonstrated under six attack families. Cross-domain transfer to speech and healthcare confirms domain independence.

\noindent\textbf{10. RobustIDPS.ai Platform.}
The complete framework is implemented as an open-source platform (DOI: 10.5281/zenodo.19129512) comprising 51 pages, 10 functional groups, and 17 integrated models. The 4-layer LLM defence framework achieves 94.5\% detection across 517 attack scenarios with four commercial LLM backends. The Multi-Agent PQC-IDS module extends detection to post-quantum protocol environments. User study validation confirms 34\% reduction in analyst triage time. The platform operates as a plug-in for Snort, Suricata, and similar IDPS frameworks.


%%---------------------------------------------------------------------
\section*{Confirmation of the Condition--Requirement Matrix}
%%---------------------------------------------------------------------

The experimental programme confirms that all twelve intersections of the $3\times4$ condition--requirement matrix are addressed with verified performance:

\begin{itemize}
\item Adversarial $\times$ Robustness: CT-TGNN Lipschitz certificates (98.3\% accuracy)
\item Adversarial $\times$ Accuracy: Stochastic Transformer calibration (ECE 0.078)
\item Adversarial $\times$ Efficiency: MambaShield $\mathcal{O}(L\log L)$ inference (0.5\,ms/flow)
\item Adversarial $\times$ Privacy: Game-theoretic certification with DP composition
\item Non-stationary $\times$ Robustness: SDE-TGNN certified robustness (76.2\%)
\item Non-stationary $\times$ Accuracy: TripleE-TGNN multi-level embeddings (96.8\%)
\item Non-stationary $\times$ Efficiency: MambaShield progressive distillation (91.4\% under poisoning)
\item Non-stationary $\times$ Privacy: Forward-compatible PQC detection (96.2\%)
\item Distributed $\times$ Robustness: FedLLM-API Byzantine resilience (87.1\% at 40\% corruption)
\item Distributed $\times$ Accuracy: FedLLM-API + LLM zero-shot (87.6\% F1)
\item Distributed $\times$ Efficiency: OT-accelerated convergence (35\% faster)
\item Distributed $\times$ Privacy: $(\varepsilon{=}0.85,\delta{=}10^{-5})$ differential privacy
\end{itemize}


%%---------------------------------------------------------------------
\section*{Industry Validation}
%%---------------------------------------------------------------------

Two implementation acts have been obtained confirming the practical applicability of the developed adversarial resilient models based on artificial intelligence for Hybrid~IDPS in network-security tasks: from \emph{Research Institute of Information and Communication Technologies LLC} (NII~IKT~LLC, Novosibirsk; integration of poisoning-resilient anomaly-detection models into the operating intrusion-detection system of the organisation and active use of the RobustIDPS.ai software platform by in-house cybersecurity specialists) and from \emph{Ksyrix~LLC} (Moscow; deployment of the developed solutions within the company's network-traffic analysis and infrastructure-protection system for detection of multi-stage attacks and coordinated anomalies on the network-interaction graph, uncertainty-based identification of previously unknown anomalous patterns, and adaptation to changes in traffic statistics). The act from the \emph{Artificial Intelligence Center of NRNU MEPhI}, Moscow, confirming application of the developed AI-based Hybrid~IDPS framework to UAV onboard-system protection (Chapter~\ref{ch:uav}), is in the course of formalisation.


%%---------------------------------------------------------------------
\section*{Directions for Further Research}
%%---------------------------------------------------------------------

The results of this dissertation open several directions for further investigation.

First, the extension of the continuous-time graph dynamics to heterogeneous temporal hypergraphs would enable modelling of higher-order interactions (e.g., multi-party transactions, group communications) that pairwise edges cannot represent, while preserving the Lipschitz certification framework.

Second, the integration of the developed methods with large multimodal foundation models (vision-language-graph) would enable joint processing of network traffic graphs with system call traces, endpoint telemetry, and DNS query streams for holistic threat detection.

Third, the deployment of quantised Multi-Agent PQC-IDS models on ARM Cortex-M and RISC-V microcontrollers would extend post-quantum traffic classification to resource-constrained IoT edge devices.

Fourth, the application of formal verification techniques (SMT-based or abstract interpretation) to prove bounded safety properties of the LLM defence pipeline would strengthen the certification guarantees beyond empirical evaluation.

Fifth, the extension of the federated learning framework to fully decentralised (peer-to-peer) settings without a trusted aggregation server would eliminate the single point of failure in the current star topology.

Sixth, the adaptation of the complete framework to additional domains --- precision agriculture (spatio-temporal crop-sensor graphs), autonomous driving (dynamic sensor-fusion graphs), and financial market surveillance (temporal transaction hypergraphs) --- would confirm the domain independence claimed in the theoretical analysis.

These directions build naturally on the mathematical foundations, algorithmic infrastructure, and software platform developed in this dissertation, and collectively aim to advance the state of the art in robust, trustworthy AI for dynamic graph-structured data.
